\section{Method}
\label{sec:method}

Cross-frequency coupling is a relation between two frequency bands, and a
token grid built from per-band content has no cell in which to write a
relation. CroFreMo adds one. It maps a multichannel EEG window to a
three-axis token grid of electrodes $\times$ frequency bands $\times$ time
patches in which every band contributes two token rows: one carrying the
band waveform, and one carrying an amplitude-preserving \emph{interaction}
representation derived from the band's measured phase coupling to slower
bands. The grid is then processed by one of two encoders: a tri-axial transformer that attends along each axis in turn, or --- the final model --- an axis-free transformer with no frequency sub-layer, in which cross-frequency relations can enter only through the token.
Two principles govern the design: coupling enters through the phase of the
interaction token and never rescales the band's amplitude features; and
every coupling pathway is gated, so that the coupling content can be
switched off within one parameterization and contributes only when
supported by optimization. The axis-free model has
2.7M parameters, the tri-axial one 1.6M.

\subsection{Cross-frequency coupling as a signal phenomenon}
\label{sec:phenomenon}

Let $x \in \mathbb{R}^{C \times T}$ be a window of $C$ electrodes at
sampling rate $f_s$. For a band-limited component $x_b(t)$, the analytic
signal $z_b = x_b + i\,\mathcal{H}[x_b]$ splits the band into an amplitude
envelope $A_b(t) = |z_b(t)|$ and a unit phase $p_b(t) = z_b(t)/|z_b(t)|
= e^{i\phi_b(t)}$. Phase--amplitude coupling (PAC) between a slow band $i$
and a fast band $j$ is the statement that $A_j(t)$ is not independent of
$\phi_i(t)$: the fast envelope is systematically larger at some phases of
the slow rhythm than at others. Its classical estimators are all functionals
of the joint distribution of $(\phi_i, A_j)$. The modulation index of
\citet{tort2010measuring} bins $\phi_i$, averages $A_j$ within each bin, and
measures the Kullback--Leibler divergence of the resulting distribution
from uniform; the mean-vector length of \citet{canolty2006high} is the
first circular moment,
\begin{equation}
\mathrm{MVL}_{ij} \;=\; \Bigl|\,\frac{1}{T}\sum_{t} A_j(t)\,e^{i\phi_i(t)}\Bigr| ,
\label{eq:mvl-classical}
\end{equation}
whose magnitude is zero when $A_j$ carries no phase information and grows
with the depth of modulation, and whose argument is the preferred phase.
Both are computed over whole recordings, with band pairs fixed in advance,
as a single scalar per pair.

CroFreMo keeps the first circular moment but changes three things about how
it is computed. A learnable sinc filterbank \citep{ravanelli2018sincnet}
with $n_b$ band-pass filters, parameterized by cutoff pairs initialized on
a logarithmic (constant-$Q$) grid, produces the band signals $x_b(t)$, so
the bands are learned rather than fixed. Time is cut into $P$
non-overlapping patches of length $\ell$ (1\,s throughout), and the moment is
taken within each patch and channel rather than over the recording, so
that transient coupling events survive at patch resolution. And the
amplitude is centered before the moment is taken,
\begin{equation}
Z_{ij} \;=\; \frac{1}{\ell}\sum_{t \in \text{patch}}
p_i(t)\,\bigl(A_j(t) - \bar A_j\bigr),
\label{eq:mvl}
\end{equation}
which removes the term $\bar A_j \cdot \frac{1}{\ell}\sum_t p_i(t)$ that a
non-uniform phase distribution alone would contribute, so that $Z_{ij}$
responds to the \emph{co-variation} of envelope and phase and not to
the mean envelope. $Z_{ij}$ is kept complex: $|Z_{ij}|$ is the strength of
band-$i$ phase coupling onto band-$j$ amplitude, $\angle Z_{ij}$ its
preferred phase. For $n_b$ bands this yields a lower-triangular
$n_b \times n_b$ complex matrix per (channel, patch), computed by the
tokenizer for every window it sees. Two caveats on what $Z_{ij}$ is.
Measured over a 1\,s patch, it contains as many cycles of band $i$ as
that band's center frequency allows, and for the slowest learned band
this can be a single cycle; and, like every phase--amplitude statistic,
it responds to co-varying non-stationarity as well as to genuine
cross-frequency interaction \citep{aru2015untangling}. We therefore treat
$Z_{ij}$ as a patch-resolution coupling statistic that the encoder may
use, not as a physiological measurement, and the results are read
accordingly.

\subsection{Coupling as token content: the duplex tokenizer}
\label{sec:tokenizer}

Each (channel, band, patch) yields three linear projections: a waveform
token $\mathbf{r}_b = W_r\,x_b[\text{patch}] \in \mathbb{R}^d$, an
amplitude feature
$\mathbf{a}_b = L_a(A_b[\text{patch}]) \in \mathbb{R}^{d/2}$, and a complex
phase feature
$\tilde{\mathbf{p}}_b = L_p(p_b[\text{patch}]) \in \mathbb{C}^{d/2}$, where
$L_p$ is linear and bias-free. The waveform token is what a per-band
tokenizer would produce on its own; the two analytic features are the raw
material for the interaction token, and the three properties below say how
that token must be built so that it carries coupling and nothing else.

\begin{proposition}[Phase-reference equivariance]
\label{prop:equivariance}
Let band $i$'s phase convention be shifted by a constant $\delta_i$, i.e.\
$p_i \mapsto e^{i\delta_i} p_i$. Then $\tilde{\mathbf{p}}_i \mapsto
e^{i\delta_i}\tilde{\mathbf{p}}_i$ and $Z_{ij} \mapsto e^{i\delta_i} Z_{ij}$
for every $j$.
\end{proposition}
\begin{proof}
$L_p$ is $\mathbb{C}$-linear without bias, so $L_p(e^{i\delta}p) =
e^{i\delta}L_p(p)$; and $e^{i\delta_i}$ factors out of the sum in
Eq.~(\ref{eq:mvl}).
\end{proof}

Raw phase depends on an arbitrary reference (the Hilbert transform fixes
it only up to the analysis convention), so by
Proposition~\ref{prop:equivariance} the phase feature
$\tilde{\mathbf{p}}_b$ is not, by itself, a well-defined token content.
We therefore let coupling enter through the \emph{aligned phase}, which
combines slower bands' phase features under the measured coupling geometry:
\begin{equation}
\mathbf{u}_j \;=\; \sum_{i<j} \alpha_{ij}\, e^{-i \angle Z_{ij}}\,
\tilde{\mathbf{p}}_i,
\qquad
\alpha_{ij} = \frac{|Z_{ij}|}{\sum_{i'<j}|Z_{i'j}|},
\label{eq:aligned}
\end{equation}
with $\mathbf{u}_0 = \tilde{\mathbf{p}}_0$ for the slowest band.

\begin{proposition}[Phase-reference invariance]
\label{prop:invariance}
For every band $j > 0$, $\mathbf{u}_j$ is invariant to the phase convention
of every band.
\end{proposition}
\begin{proof}
Under $p_i \mapsto e^{i\delta_i} p_i$, Proposition~\ref{prop:equivariance}
gives $\tilde{\mathbf{p}}_i \mapsto e^{i\delta_i}\tilde{\mathbf{p}}_i$ and
$e^{-i\angle Z_{ij}} \mapsto e^{-i\delta_i}e^{-i\angle Z_{ij}}$; the two
factors cancel term by term in Eq.~(\ref{eq:aligned}), and $\alpha_{ij}$
depends on $|Z_{ij}|$ only.
\end{proof}

The slowest band has no slower band to align to, so $\mathbf{u}_0 =
\tilde{\mathbf{p}}_0$ is its own phase feature and is not
reference-invariant; its interaction token carries analytic-signal
content but no coupling geometry, and the claim of
Proposition~\ref{prop:invariance} excludes it.
Proposition~\ref{prop:invariance} is the reason coupling enters through
Eq.~(\ref{eq:aligned}) rather than through the raw phases: the rotation by
$e^{-i\angle Z_{ij}}$ re-expresses each slow band's phase feature relative
to the phase at which it actually modulates band $j$, which is a physical
quantity, and the weights $\alpha_{ij}$ let the strongest modulators of
band $j$ dominate its aligned phase.

\begin{proposition}[Magnitude preservation]
\label{prop:isometry}
Define the interaction token
\begin{equation}
\mathbf{h}_j \;=\; \mathbf{a}_j \odot
\frac{\mathbf{u}_j}{|\mathbf{u}_j|} .
\label{eq:rotation}
\end{equation}
Then $|\mathbf{h}_j| = |\mathbf{a}_j|$ elementwise, for every value of the
coupling.
\end{proposition}
\begin{proof}
Each component of $\mathbf{u}_j/|\mathbf{u}_j|$ has unit modulus.
\end{proof}

Multiplying by $\mathbf{u}_j$ directly would rescale the amplitude feature
by $|\mathbf{u}_j|$, confounding coupling geometry with changes in band
power, which is the one quantity every EEG model already exploits. Under
Eq.~(\ref{eq:rotation}) the amplitude feature reaches the encoder with
its modulus untouched and the coupling geometry is encoded purely in
phase. Two consequences should be stated plainly. First, the weights
$\alpha_{ij}$ of Eq.~(\ref{eq:aligned}) normalize the column
$Z_{\cdot j}$, so multiplying that column by any positive constant changes
neither $\alpha_{ij}$ nor the angles: the interaction token keeps the
\emph{relative} coupling weights and the alignment geometry, not the
absolute coupling strength, and near-zero coupling does not make the
token vanish. Second, Proposition~\ref{prop:isometry} constrains the
interaction token alone; the fused row $\mathbf{r}_b + \boldsymbol{\beta}_b
\odot \mathbf{h}_b$ of Eq.~(\ref{eq:duplex}) has no such constraint, and
the interaction row hands the encoder an explicit amplitude-envelope
feature that a waveform-only grid does not expose. The ablation of
Section~\ref{sec:attribution} therefore attributes gains to the
interaction rows as a whole; separating the envelope from the
cross-frequency alignment needs the matched controls discussed there.

\paragraph{The duplex grid.}
The final grid stacks, for each physical band $b$, a \emph{fused} row and
a gated \emph{interaction} row:
\begin{equation}
\text{row}_b = \mathbf{r}_b + \boldsymbol{\beta}_b \odot \mathbf{h}_b, \qquad
\text{row}_{n_b + b} = \gamma_b\, \mathbf{h}_b,
\label{eq:duplex}
\end{equation}
with $\boldsymbol{\beta}_b \in \mathbb{R}^d$ initialized to zero and scalar
$\gamma_b$ initialized to one. Here $\mathbf{h}_b \in \mathbb{C}^{d/2}$ is
laid out as the real and imaginary parts of its $d/2$ components, giving
a real vector of dimension $d$ on which the elementwise gate
$\boldsymbol{\beta}_b$ and the addition to $\mathbf{r}_b$ act. The two rows give coupling two different
routes into the encoder. The fused row lets coupling modulate the band's
own waveform representation in place; because $\boldsymbol{\beta}_b$
starts at zero, its gradient at initialization is
$\partial\mathcal{L}/\partial\boldsymbol{\beta}_b =
\partial\mathcal{L}/\partial\,\text{row}_b \odot \mathbf{h}_b$, so
coupling enters this route only along the directions in which the
interaction token is correlated with the loss gradient, and the fused rows
coincide with the uncoupled model until it does. The interaction row makes
coupling a separately attendable feature from the first step, which is what
event-level tasks turn out to need (Section~\ref{sec:results}).

\begin{proposition}[Gating]
\label{prop:recover}
At $\boldsymbol{\beta}_b = 0$ the fused row of band $b$ equals its waveform
token $\mathbf{r}_b$ exactly. At $\gamma_b = 0$ the interaction row of
band $b$ is the zero vector before positional embedding, and carries no
coupling content.
\end{proposition}
\begin{proof}
Both statements read off Eq.~(\ref{eq:duplex}).
\end{proof}

We do not claim more than this. A zero interaction row still occupies a
position in self-attention: it contributes a term to every query's
softmax normalization and, once positional embeddings are added, is not
zero as a key or value, so the configuration $\boldsymbol{\beta} = 0$,
$\gamma = 0$ is not function-identical to the $n_b$-row waveform model;
it is the waveform model plus $n_b$ learned bias rows per electrode.
Coupling content can be switched off within the parameterization; the
row count cannot. The waveform-only comparisons in
Section~\ref{sec:results} are therefore run as separate $n_b$-row models,
and the trained values of $\boldsymbol{\beta}$ and $\gamma$ are read as
measurements of how much coupling content each task retained.

\subsection{Encoders: cross-frequency structure, or none}
\label{sec:encoder}

The grid $(C \times 2n_b \times P \times d)$ can be processed by an encoder
that supplies cross-frequency relations itself, or by one that does not.
We build both, with the same tokenizer, the same recipe and the same
mean-pooled linear head, because the difference between them is the
experiment of Section~\ref{sec:content-structure}.

\paragraph{Tri-axial encoder (structure).}
$N$ pre-norm blocks each apply self-attention along one axis at a time ---
time (with rotary position encoding), electrodes (with a learned
per-electrode embedding), and the $2n_b$ frequency rows (with an embedding
of each band's learned center frequency and bandwidth) --- followed by a
feed-forward layer, each as a residual. The frequency sub-layer is plain
attention over rows: it gives the encoder a dedicated path along which
cross-band relations can be formed from the waveform rows, so the
interaction token competes with structure that may re-derive some of what
it carries. We use
$n_b{=}8$, $d{=}128$, $N{=}6$, four heads, 1.6M parameters.

\paragraph{Axis-free encoder (content only).}
The final model removes the frequency sub-layer entirely and folds the band
rows into the spatial one: in each block, attention runs first along time
per (electrode, row), then jointly over the $C \cdot 2n_b$ tokens of each
patch, then through the feed-forward layer. Rows of different bands meet
only inside this folded spatial attention, which carries no band
embedding and no frequency-specific parameters. The encoder therefore has
no frequency-specific structure: bands can still exchange information
through that attention, and cross-band relations can in principle be
learned from waveform rows alone, but nothing in the architecture is
dedicated to them. The interaction token of Eq.~(\ref{eq:rotation}) is
the only \emph{explicit} cross-frequency representation the encoder
receives, and the waveform-only model of Section~\ref{sec:attribution}
is the same encoder without it. Width is raised
to $d{=}192$ to spend the parameters the frequency sub-layer no longer
uses; $N{=}6$, four heads, 2.7M parameters. Section~\ref{sec:design}
reports the alternatives (no folding, $d{=}256$, an explicit
coupling-strength feature) that were tried and not adopted.

\paragraph{The token inside a published encoder.}
To test the same distinction in an encoder we did not design, we add the
tokenizer's rows to CBraMod \citep{wang2025cbramod} without removing
anything: CBraMod's patch embedding (a convolutional stem plus an rFFT
magnitude branch), its criss-cross encoder, positional convolution and
classifier are kept verbatim, and each electrode receives $n_b$ extra
channels carrying either the $n_b$ interaction tokens or, as the matched
control, the $n_b$ waveform tokens of the same frontend. The two arms have
the same parameter count and the same token count and differ only in what
the extra rows carry. CBraMod's own spectral branch remains in place, so
this encoder, unlike the axis-free one, is never starved of
cross-frequency information.

\subsection{A pretraining objective for coupling}
\label{sec:objective}

If cross-frequency structure belongs in the token, it is natural to ask
whether it also belongs in the pretraining target. We construct the
objective that puts it there; Section~\ref{sec:results} evaluates it
empirically, to convergence.

\paragraph{Masking.} A random subset $\mathcal{M}$ of \emph{physical}
bands is masked per channel. Because the interaction row of band $b$
carries $\mathbf{a}_b$ by construction (Eq.~(\ref{eq:rotation})), masking
hides \emph{both} of $b$'s rows --- fused and interaction --- or the
visible interaction row would hand the reconstruction its own target; both
are replaced by a shared learned mask token before the positional
embeddings. Entries of the coupling matrix consumed by the encoder are
zeroed wherever either band is masked, so the target cannot travel through
the coupling input.

\paragraph{Objective.} The loss combines an amplitude term and a coupling
term,
\begin{equation}
\mathcal{L} \;=\; \mathcal{L}_{\mathrm{amp}}
\;+\; \lambda\, \mathcal{L}_{\mathrm{cfc}},
\qquad \lambda = 1.
\label{eq:objective}
\end{equation}
$\mathcal{L}_{\mathrm{amp}}$ asks the model to predict, from each masked
band's fused-row output, the patch log-amplitude standardized \emph{within
each band} across the batch --- the rebalancing of
\citet{yu2026understanding}, adopted as-is --- so that every band's
structure is rewarded equally rather than in proportion to its
$1/f$-dominated variance:
\begin{equation}
\mathcal{L}_{\mathrm{amp}} = \frac{1}{|\mathcal{M}|}
\sum_{j \in \mathcal{M}}
\bigl\| \hat{y}_j - \mathrm{std}_b\!\bigl(\log \bar A_j\bigr) \bigr\|^2 .
\label{eq:lamp}
\end{equation}
$\mathcal{L}_{\mathrm{cfc}}$ asks it to predict the masked band's coupling
column --- the real and imaginary parts of $Z_{ij}$ --- restricted to
\emph{visible} phase bands $i$, since the phase required to define
$Z_{ij}$ is unavailable when band $i$ is masked:
\begin{equation}
\mathcal{L}_{\mathrm{cfc}} = \frac{1}{Z_{\mathrm{sup}}}
\sum_{j \in \mathcal{M}} \sum_{i \notin \mathcal{M}}
\bigl\| \hat{Z}_{ij} - Z_{ij} \bigr\|^2,
\label{eq:lcfc}
\end{equation}
with $Z_{\mathrm{sup}}$ the count of supervised entries. Pretraining uses
the same corpora pool for every objective variant compared in
Section~\ref{sec:results}; per-variant schedules are given in the appendix.
